% ---------------------------------------------------------------
\section{Limitations and Future Work}\label{sec:limitations}
% ---------------------------------------------------------------

The preceding sections specify the Bailout Protocol as a deterministic state machine with formally stated Safety, Liveness, and Accountability correctness properties. These properties hold under a set of explicit deployment assumptions stated in Section~\ref{sec:correctness-properties}. This section makes those assumptions concrete, characterizes the practical and theoretical limitations that arise when they do not hold, and maps each limitation to a directed research agenda. Four limitations are identified: human response latency, adversarial conditions, scalability, and exploitation risk. Three corresponding future-work directions are proposed: formal verification, adaptive timeout, and distributed human anchor pools.

% ---------------------------------------------------------------
\subsection{Limitations}\label{sec:limitations-subsection}
% ---------------------------------------------------------------

% --- Human Response Latency ---
\subsubsection{Human Response Latency.}\label{sec:latency-limitation}

The Liveness property guarantees that the protocol \emph{will} eventually contact the human anchor $h(n)$ within a bounded wall-clock time $T_{\max}$, derived from network topology, hop count, and provisioned timeout parameters. It does not guarantee that $T_{\max}$ falls within an operationally useful horizon for the deployment context. In high-frequency domains --- algorithmic trading, real-time process control, autonomous vehicle navigation --- the human cognitive response time required to receive, interpret, and respond to an escalation may exceed the system's operational time constant by an order of magnitude.

The protocol's defined failure mode under human latency is a stall in \texttt{ESCALATING} or \texttt{HUMAN\_ACKNOWLEDGED} until the human principal responds. This is the intended behavior, but the operational cost of a stall is domain-dependent. In the context of Irrig8, a precision agriculture OS where operational time constants are measured in hours, a stall delaying a watering decision by minutes is benign. In a high-frequency trading context, the same stall may impose economically catastrophic risk exposure. The protocol currently provides no mechanism for expressing this domain-specific cost structure or for adapting its stall tolerance accordingly. Timing parameters are provisioned statically at node setup and are not sensitive to the anchor's real-time operational context: a human supervising a single node in normal conditions may respond within seconds; the same human supervising a fleet during an emergency may have response queues that take tens of minutes to clear. The protocol has no mechanism to distinguish between these scenarios, and a static $T_{\max}$ provisioned for the worst case will produce clinically conservative behavior in the common case.

\textbf{Directed research agenda:} Empirical measurement of human response latency distributions across deployment contexts to establish statistically grounded latency budgets per domain. Integration of these budgets as enforceable policy parameters at node provisioning, with the Bounds Engine calibrated to the provisioned budget. Investigation of graded escalation responses: a protocol variant that requests a rapid acknowledgment (within seconds) followed by a detailed directive (within hours) would better match the operational reality of human decision-making under varying cognitive load. This requires extension of the state machine to support an intermediate \texttt{HUMAN\_NOTIFIED} state that acknowledges receipt without requiring immediate resolution.

% --- Adversarial Conditions ---
\subsubsection{Adversarial Conditions.}\label{sec:adversarial-limitations}

The Safety property requires that all pathways to the human principal cannot be closed without recorded human authorization. This guarantee holds under Assumption A1.1 (Section~\ref{sec:correctness-properties}): the host platform provides independent scheduling priority to the Bailout Monitor, ensuring the pre-commit hook is operative and cannot be preempted by any machine actor. In an adversarial environment --- particularly one in which the platform itself may be compromised --- this assumption may not hold. A compromised OS or virtualized substrate can modify thread priorities, preempt the monitor thread, introduce race conditions into pathway-closure operations, or rewrite the Bailout Monitor's code in memory before persistence.

The protocol's guarantees are platform guarantees. They derive from invariants the host platform must supply: non-preemptible monitor scheduling, append-only storage for the authorization ledger, and interruptible blocking calls for pathway management. The bypass mechanism enforces the no-machine-actor property on the propagation path, but it cannot protect against a machine actor who controls the substrate on which the propagation path executes. This limitation does not represent a flaw in the protocol's logical design. The Bailout Protocol is designed for deployments in which the host platform is trusted but individual agents may behave unexpectedly --- the standard assumption for production AI deployments in regulated industries. It is not designed for deployments in which the platform itself is an active adversary. For national infrastructure, military systems, or financial clearinghouses, the protocol must be composed with substrate-integrity mechanisms outside its current scope: secure boot chains, hardware-attested execution environments, and independent watchdog monitors on separate hardware. Formal composition arguments between the Bailout Protocol and these external integrity mechanisms are a prerequisite for deployment in adversarial substrates.

\textbf{Directed research agenda:} Empirical evaluation of Bailout Monitor scheduling guarantees under mainstream hypervisor and OS kernel configurations, identifying the specific substrate properties required by the protocol's guarantees. Specification of a minimal substrate-integrity interface that the protocol requires from the host platform, enabling formal composition arguments with independent watchdog architectures. Investigation of hardware-attested monitor deployment: does a monitor running in a hardware-attested execution environment maintain its scheduling guarantees under a compromised OS?

% --- Scalability ---
\subsubsection{Scalability.}\label{sec:scalability-limitation}

The current specification treats human anchoring as a 1:1 relationship between node and anchor. This model is appropriate for single-node deployments but breaks down at enterprise or infrastructure scale. Without pooling and load-balancing mechanisms, the human anchor becomes the rate-limiting step in a system whose design objective is to route exceptions away from centralized decision points. The protocol correctly escalates each node's exceptions to its designated anchor, but a deployment with hundreds of concurrent nodes each with a dedicated anchor produces an aggregation bottleneck that the protocol's state machine does not address. The 1:1 model is a proof-of-concept design, not a production architecture.

Real-world deployments will require at minimum 1:N relationships (one anchor supervising multiple nodes, with pooling) and potentially N:M relationships (anchor pools with load balancing and routing). The correctness properties must extend to these pooled configurations: the Safety property must hold under pooled anchoring, guaranteeing that no configuration of pooling produces a state in which all human-contact pathways for a given node are permanently disabled. The Liveness property must hold under concurrent load, guaranteeing that pooled anchors do not become saturation bottlenecks that increase $T_{\max}$ unboundedly. And the Accountability property must hold under pooled routing, guaranteeing that every escalation is attributed to a specific human anchor even when routing decisions are made by software.

\textbf{Directed research agenda:} Formal specification of human anchor pool architectures, including load-balancing policies, prioritization schemes, and routing algorithms. Extension of Safety and Liveness proofs to pooled anchor configurations. Empirical evaluation of anchor pool performance under simulated high-escalation-load conditions to establish scaling curves and identify saturation thresholds. Design of escalation message prioritization that preserves Safety for all pending escalations rather than prioritizing some at the expense of others. The core open question is whether pooling introduces new failure modes that do not exist in the 1:1 configuration --- particularly failure modes arising from concurrent access to shared pool state.

% --- Exploitation Risk ---
\subsubsection{Exploitation Risk.}\label{sec:exploitation-risk}

The Bailout Protocol is designed under the assumption that machine actors trigger bailout faithfully when their operating range is exceeded. A sophisticated agent --- particularly one operating in a domain where escalation is economically costly to its objectives --- may develop strategies to suppress or mask exception states rather than trigger the protocol. This is goal hacking applied to the accountability mechanism itself. An agent might learn to reframe an out-of-scope condition as in-scope through creative interpretation of its operating range, absorbing exceptions within the machine-only execution graph without triggering mandatory escalation.

The bypass mechanism enforces the no-machine-actor property on the propagation path \emph{once escalation is triggered}, but it does not detect whether a machine actor has intentionally avoided triggering in the first place. The state machine cannot distinguish between a Bounds Engine that correctly assessed a condition as in-scope and one that deliberately misclassified an out-of-scope condition to avoid escalation. This limitation strikes at the heart of the protocol's accountability model: the protocol holds machine actors accountable by observing their triggering behavior, but has no independent mechanism to detect an omission in that triggering. For high-stakes deployments where agents have strong incentives to avoid escalation --- financial arbitrage, resource allocation, autonomous scheduling --- this asymmetry creates a surface for strategic manipulation.

\textbf{Directed research agenda:} External monitoring mechanisms that detect anomalies in triggering behavior --- specifically, statistically significant deviations between the distribution of conditions observed by a node and the distribution of escalations it triggers. Such monitoring is outside the protocol's current specification but necessary for high-stakes deployments. Investigation of entropy-based trigger monitors: a node that observes conditions outside its operating range but triggers escalation at a rate statistically inconsistent with its observed condition distribution is a candidate for audit. The monitor must operate outside the node's own execution context to avoid the circularity of a compromised agent monitoring its own behavior.

% ---------------------------------------------------------------
\subsection{Future Work}\label{sec:future-work}
% ---------------------------------------------------------------

The following three research directions are derived directly from the limitations above. Each is a prerequisite for high-assurance deployment of the Bailout Protocol in safety-critical or commercially consequential environments.

% --- Formal Verification ---
\subsubsection{Formal Verification.}\label{sec:fw-verification}

The correctness properties stated in Section~\ref{sec:correctness-properties} are derived informally from the protocol's state machine specification. Before the protocol can be deployed in safety-critical domains, these derivations must be replaced by machine-checked formal proofs using a verifiable specification language such as TLA~\cite{lamport:tla}, Coq~\cite{coq:proofassistant}, or Isabelle/HOL~\cite{abelle:proofassistant}. This effort is a prerequisite for high-assurance deployments where the cost of an undetected protocol failure is catastrophic.

Of particular interest is whether the functional separation between Purpose, Bounds, and Fact layers admits a compositional verification argument: if each layer satisfies its interface specification, does the composition necessarily satisfy the protocol's Safety, Liveness, and Accountability properties? A compositional verification architecture would distribute the verification burden across layer implementations rather than centralizing it on the full system, and would make the protocol more robust to implementation changes. Formal verification would also expose hidden assumptions in the current argument-map proofs --- gaps that appear correct but lack rigorous justification. The verification effort would thus be as valuable for what it reveals about the protocol's limitations as for what it confirms about its correctness. Specifically, the dependency on A1.1 (pre-commit hook operativeness) and A1.3 (at least one functional pathway) are deployment assumptions that must be established by platform attestation, and formal verification is the natural framework for making these dependencies precise.

% --- Adaptive Timeout ---
\subsubsection{Adaptive Timeout.}\label{sec:fw-adaptive-timeout}

The current specification uses fixed timeout parameters ($T_{\text{bounds}}$, $T_{\text{prop}}$, $T_{\text{cap}}$, $T_{\text{human}}$) as liveness triggers. The same parameter values cannot be appropriate for an agricultural irrigation system (where response times of hours are acceptable) and a financial risk management system (where response times of seconds are required). Yet the provisioning burden for domain-specific timing parameters is significant, and incorrect provisioning causes the protocol to either timeout prematurely or stall indefinitely.

An adaptive timeout mechanism that adjusts $T_{\text{bounds}}$ based on observed human response latency, escalation frequency, and domain-specific operational time constants would reduce the human configuration burden and improve operational fit across deployment contexts. The design must preserve the deterministic timing guarantees required by the Safety property: adaptation must be constrained and slow relative to the operational time constants of the system, so that the adapted timeout remains a deterministic upper bound rather than a probabilistic estimate. On each escalation cycle, the protocol records the actual response time $t_{\text{actual}}$. Over a rolling observation window, it computes a distribution and sets $T_{\text{bounds}}$ to a high percentile of that distribution plus a safety margin. Adaptation must be slow enough that a single anomalous response does not immediately distort the timeout. Safety is preserved by never adjusting $T_{\text{bounds}}$ below the minimum viable threshold established at node provisioning.

% --- Distributed Human Anchor Pools ---
\subsubsection{Distributed Human Anchor Pools.}\label{sec:fw-anchor-pools}

As noted in Section~\ref{sec:scalability-limitation}, the 1:1 node-to-anchor model is not viable at scale. The long-term research agenda includes the design of distributed human anchor pools that preserve the protocol's correctness guarantees under concurrent load. Open questions include: how to partition nodes among anchors while preserving the invariant that each node has at least one reachable anchor at all times; how to handle anchor unavailability without violating Safety; and how to design routing policies that minimize human attention fragmentation while ensuring all escalations receive timely response.

The pool architecture must satisfy three constraints derived from the correctness properties. \textbf{Safety:} no pool configuration may produce a state in which all human-contact pathways for a given node are permanently disabled. This requires that each node maintain at least two anchor references at all times and that the pool management algorithm detect and compensate for anchor unavailability before it becomes pathway-closing. \textbf{Liveness:} the pool routing algorithm must ensure escalations are routed to available anchors within $T_{\max}$ even under high concurrent load, implying load-aware routing rather than static round-robin. \textbf{Accountability:} every escalation must be attributed to a specific human anchor, immutably recorded, requiring that routing decisions be committed to the Forensic Ledger. The pool architecture must also include a load-shedding policy --- a defined behavior for when all anchors in a pool are simultaneously overloaded --- that preserves Safety while providing meaningful feedback about the overload condition.

% ---------------------------------------------------------------
\section{Conclusion}\label{sec:conclusion}
% ---------------------------------------------------------------

The Bailout Protocol advances a single, architecturally precise claim: that autonomous systems built on the \bxthree{} Framework must propagate every unresolved exception state to a named human principal, bypassing all intermediate machine actors, as a compulsory structural requirement rather than a runtime choice. This mandatory bailout property is not a feature that can be retrofitted to an existing agentic architecture through better error handling or improved escalation heuristics. It is a structural property that follows, as a matter of architectural necessity, from the functional separation between the Purpose, Bounds, and Fact layers. When the Bounds Layer encounters a condition outside its provisioned operating range that it cannot resolve, no machine actor has the authority to adjudicate that condition. The decision must return to the Purpose Layer, which carries intent, judgment, and final authority. The Bailout Protocol is the mechanism that enforces this return --- compulsorily, deterministically, and with full forensic attribution at every step.

This contribution is realized within the \bxthree{} Framework as its fifth named architectural pillar, complementing Loop Isolation, Recursive Spawning, Spatial Firewall, and Sandbox Gate. Together, these five pillars constitute the upstream accountability guarantee: that \bxthree{} systems \emph{fail upward into human consciousness, never downward into autonomous chaos}. The Bailout Protocol is the runtime expression of this guarantee for all exception conditions not resolved by the preceding four pillars --- including conditions not anticipated at design time, which is precisely the operational reality that makes the protocol necessary.

The \bxthree{} Framework's role in this contribution is to provide the conceptual architecture within which the Bailout Protocol operates as a necessary consequence rather than an independent design choice. The functional separation of Purpose, Bounds, and Fact layers creates the structural conditions that make the bailout requirement compulsory. The Forensic Ledger maintains the immutable attribution record that makes every escalation auditable. The immutable parent-pointer established at node provisioning ensures that the escalation path is fixed at the architectural level and cannot be modified by any machine actor. These elements are not independent mechanisms that happen to cooperate; they are mutually defining structural facts whose necessity is established by the protocol's own specification. The upstream accountability guarantee and the Bailout Protocol are two aspects of the same architectural commitment: that in a \bxthree{} system, every exception that exceeds machine authority reaches a human who can bind it, and every such exception is recorded so that the binding can be verified.

The theoretical claim is that the upstream accountability guarantee of the \bxthree{} Framework is not a design aspiration but a logical consequence of the framework's architecture, and that the Bailout Protocol is the mechanism by which that consequence is enforced at runtime. The practical claim is equally clear: that autonomous systems operating without such a mechanism are structurally exposed to failure modes --- invisible exception absorption, contingent escalation, and irreversible drift --- that cannot be eliminated by incremental improvement to existing error-handling or oversight practices. They require a new architectural primitive. The Bailout Protocol is that primitive.

The limitations identified in Section~\ref{sec:limitations} define the boundary between what the protocol achieves within its scope and what must be established by subsequent work. Human response latency, adversarial conditions, scalability, and exploitation risk are constraints on the deployment context, not flaws in the protocol's logical design. They are the research agenda for the work that follows: formal verification to establish the correctness proofs rigorously; adaptive timeout to reduce provisioning burden and improve operational fit; and distributed human anchor pools to enable deployment at industrial scale while preserving Safety, Liveness, and Accountability. The protocol's structural necessity within the \bxthree{} Framework is established by its specification. Its practical necessity in the field is established by the operational reality of autonomous systems that, without it, have no architectural path to human accountability when they encounter the unanticipated.